\section{Results}

\subsection{H1 rejected: coherence does not causally improve faithfulness at fixed relevance}

The matched construction succeeds: 200 SQuAD query pairs are built with mean
absolute similarity gap 0.006 and mean CCS gap 0.533. This is a strong
manipulation of set coherence while local relevance is nearly fixed. The
faithfulness effect is null (Table~\ref{tab:causal}).

\begin{table}[H]
\centering
\small
\caption{Matched-similarity intervention. Higher CCS does not improve
faithfulness at fixed mean query similarity.}
\label{tab:causal}
\begin{tabular}{lrrrrrr}
\toprule
Setting & $n$ & HIGH faith & LOW faith & HIGH--LOW & $p$ & 95\% CI \\
\midrule
SQuAD / Mistral / DeBERTa & 200 & 0.636 & 0.639 & -0.002 & 0.628 & [-0.022, 0.017] \\
\bottomrule
\end{tabular}
\end{table}

The binary hallucination result is even less compatible with a simple causal
story: HIGH-CCS sets hallucinate more often than LOW-CCS sets (16.5\% vs.
9.0\%). We therefore reject H1. The result turns the old coherence story into
evidence for the paper's current thesis: without a matched control, a plausible
retrieval-set statistic can be mistaken for a faithfulness mechanism.

\begin{figure}[H]
\centering
\includegraphics[width=0.78\linewidth]{fix_01_paired_diff.pdf}
\caption{Paired HIGH--LOW CCS faithfulness differences under the
matched-similarity intervention. The distribution is centered near zero, and
the Wilcoxon test does not reject equality ($p=0.628$), so the controlled test
rejects the causal coherence hypothesis rather than supporting it.}
\label{fig:matched_ccs_null}
\end{figure}

\subsection{H2 weakened: the scaled DeBERTa refinement effect is small}

Scaling the SQuAD/Mistral cell to $n=500$ over five seeds reduces the
DeBERTa-measured baseline-minus-refinement effect to 1.1 faithfulness points
(Table~\ref{tab:scaled}). Only one of five seeds is significant for the
baseline-minus-HCPC-v1 contrast. HCPC-v2 restores the HCPC-v1 drop on average,
but the effect is not the large, stable phenomenon suggested by the pilot.

\begin{table}[H]
\centering
\small
\caption{Scaled SQuAD/Mistral headline cell.}
\label{tab:scaled}
\begin{tabular}{lrrrr}
\toprule
Condition & $n$ & Faithfulness & Hallucination & Retrieval similarity \\
\midrule
Baseline & 2500 & 0.661 & 0.146 & 0.532 \\
HCPC-v1 & 2500 & 0.650 & 0.147 & 0.569 \\
HCPC-v2 & 2500 & 0.661 & 0.143 & 0.532 \\
\bottomrule
\end{tabular}
\end{table}

This is still an interesting non-monotonicity: HCPC-v1 raises retrieval
similarity while lowering DeBERTa faithfulness. But it is a small effect under
this metric, not a headline-scale failure. The original $n=30$ result should
therefore be treated as pilot evidence.

\subsection{H3 rejected: the same generations imply different conclusions}

Metric choice changes the magnitude of the core result
(Table~\ref{tab:metrics}). This is the most important positive result in the
revision: evaluator choice is not a reporting detail, but part of the
identified estimand. The same 7,500 generations yield a 1.1 point
baseline-minus-HCPC-v1 effect under DeBERTa, 3.2 points under a second NLI
model, and 14.0 points under a local RAGAS-style judge. DeBERTa and RAGAS
correlate only $r=0.182$; the second NLI model and RAGAS correlate $r=0.674$.

\begin{table}[H]
\centering
\small
\setlength{\tabcolsep}{7pt}
\caption{Central metric-dependence result: identical generations produce
different intervention magnitudes under different faithfulness judges.}
\label{tab:metrics}
\begin{tabular}{lrrrr}
\toprule
Condition & $n$ & DeBERTa & Second NLI & RAGAS-style judge \\
\midrule
Baseline & 2500 & 0.661 & 0.350 & 0.730 \\
HCPC-v1 & 2500 & 0.650 & 0.318 & 0.590 \\
HCPC-v2 & 2500 & 0.661 & 0.351 & 0.728 \\
\midrule
Baseline--v1 & -- & 0.011 & 0.032 & 0.140 \\
\bottomrule
\end{tabular}
\end{table}

\begin{table}[H]
\centering
\small
\caption{Metric agreement on the same 7,500 generations.}
\label{tab:metric_corr}
\begin{tabular}{lrr}
\toprule
Metric pair & Pearson $r$ & Spearman $\rho$ \\
\midrule
DeBERTa / second NLI & 0.259 & 0.265 \\
DeBERTa / RAGAS-style judge & 0.182 & 0.212 \\
Second NLI / RAGAS-style judge & 0.674 & 0.651 \\
\bottomrule
\end{tabular}
\end{table}

A single faithfulness metric does not identify the effect size of a RAG
intervention. The result is not merely that metrics have different scales; the
weak DeBERTa--RAGAS correlation means that the evaluator changes the empirical
conclusion.

\subsection{H4 dataset-dependent: threshold transfer is uneven}

For evaluation dataset $d$ and threshold $\tau$, the recovery score is
\[
\mathrm{Recovery}(d,\tau) =
\frac{\mathrm{Faith}_{\mathrm{gate}}(d,\tau) -
\mathrm{Faith}_{\mathrm{HCPC}\mbox{-}\mathrm{v1}}(d)}
{\mathrm{Faith}_{\mathrm{baseline}}(d) -
\mathrm{Faith}_{\mathrm{HCPC}\mbox{-}\mathrm{v1}}(d)}.
\]
The diagonal recovery is the recovery obtained when $\tau$ is tuned and
evaluated on the same dataset. The off-diagonal mean averages recovery when
that tuned $\tau$ is transferred to the other four datasets. The gap is their
difference. Scores can exceed 1 when the gated policy exceeds the baseline
relative to HCPC-v1, and they can be negative when the transferred threshold is
worse than HCPC-v1 under this normalization.

The $5\times5$ threshold-transfer matrix shows large in-domain advantages for
SQuAD, PubMedQA, and NaturalQuestions (Table~\ref{tab:tauflags}). HotpotQA and
TriviaQA do not show the same pattern. Thus, a single threshold tuned on one
dataset cannot be treated as generally validated.

\begin{table}[H]
\centering
\small
\caption{Threshold-transfer recovery score: in-domain versus cross-dataset
threshold effect.}
\label{tab:tauflags}
\begin{tabular}{lrrr}
\toprule
Tune dataset & Diagonal recovery & Off-diagonal mean & Gap \\
\midrule
PubMedQA & 1.452 & 0.451 & 1.002 \\
NaturalQuestions & 1.401 & 0.211 & 1.191 \\
SQuAD & 0.823 & -0.141 & 0.963 \\
TriviaQA & 0.464 & 0.698 & -0.234 \\
HotpotQA & -0.112 & 0.093 & -0.205 \\
\bottomrule
\end{tabular}
\end{table}

\subsection{Coherence remains diagnostically useful}

The noise experiment prevents an overcorrection. Although CCS is not causal in
Test 1, coherence is not equivalent to generic retrieval degradation. Random
off-topic noise has a steeper faithfulness slope than same-topic
answer-uninformative noise: $-0.069$ vs. $-0.043$ per unit noise rate
(Table~\ref{tab:noise}). Same-topic noise still hurts, but it hurts less.

\begin{table}[H]
\centering
\small
\caption{Noise slope response.}
\label{tab:noise}
\begin{tabular}{lrrr}
\toprule
Condition & Faith slope & Similarity slope & Drop at full noise \\
\midrule
Random noise & -0.069 & -0.481 & 0.052 \\
Coherent uninformative noise & -0.043 & -0.113 & 0.042 \\
HCPC-v1 refinement & -- & -- & 0.001 \\
\bottomrule
\end{tabular}
\end{table}

This supports a narrower research claim: set-level context structure is a
useful diagnostic variable, but it is neither sufficient for support nor
validated as an empirical mechanism.

\subsection{H5 rejected: no policy is uniformly best after cost is included}

Figure~\ref{fig:pareto_p99} and Table~\ref{tab:h2h} report the full
SQuAD/HotpotQA baseline comparison using the verified Self-RAG artifact.
HCPC-v2 is competitive on SQuAD, but RAPTOR-2L is slightly higher in
faithfulness and lower in p99 latency. On HotpotQA, CRAG wins on faithfulness,
hallucination rate, and latency. Self-RAG performs poorly in this matched
short-answer harness.

\begin{figure}[H]
\centering
\includegraphics[width=0.98\linewidth]{fix_06_pareto_p99_log_full_selfrag.pdf}
\caption{Cost-aware Pareto view of the full baseline audit. The x-axis uses
p99 latency on a log scale so the Self-RAG tail remains visible without
compressing the other methods. No method is uniformly best across datasets once
faithfulness, tail latency, hallucination behavior, and offline indexing costs
are considered.}
\label{fig:pareto_p99}
\end{figure}

\begin{table}[H]
\centering
\scriptsize
\caption{Head-to-head baseline comparison, $n=200$ per dataset/condition.
Latency is milliseconds.}
\label{tab:h2h}
\begin{tabular}{llrrrr}
\toprule
Dataset & Method & Faith & Halluc. & Mean lat. & p99 lat. \\
\midrule
SQuAD & CRAG & 0.698 & 0.120 & 1021 & 3318 \\
SQuAD & HCPC-v2 & 0.708 & 0.125 & 1349 & 3612 \\
SQuAD & RAPTOR-2L & 0.710 & 0.125 & 1296 & 3189 \\
SQuAD & Self-RAG & 0.574 & 0.390 & 7542 & 44800 \\
\midrule
HotpotQA & CRAG & 0.643 & 0.105 & 1473 & 3281 \\
HotpotQA & HCPC-v2 & 0.632 & 0.130 & 1841 & 4032 \\
HotpotQA & RAPTOR-2L & 0.631 & 0.130 & 1906 & 4043 \\
HotpotQA & Self-RAG & 0.574 & 0.430 & 6158 & 46965 \\
\bottomrule
\end{tabular}
\end{table}

RAPTOR's offline tree-building cost is also nontrivial: 99.7--161.2 seconds
across datasets, compared with 1.2--8.3 seconds for dense indexing
(Table~\ref{tab:raptorcost}). Method claims that omit indexing and tail latency
miss part of the deployment result.

\begin{table}[H]
\centering
\scriptsize
\caption{RAPTOR cost audit across three datasets, $n=100$ each.}
\label{tab:raptorcost}
\begin{tabular}{lrrrrrr}
\toprule
Dataset & Faith & Halluc. & p50 lat. & p99 lat. & Dense idx. & RAPTOR idx. \\
\midrule
SQuAD & 0.789 & 0.050 & 1190 & 4943 & 8.333 s & 161.203 s \\
PubMedQA & 0.560 & 0.290 & 3901 & 6812 & 1.177 s & 108.264 s \\
HotpotQA & 0.617 & 0.210 & 1973 & 4274 & 7.140 s & 99.740 s \\
\bottomrule
\end{tabular}
\end{table}
